
\section{Estudio adversarial de identificabilidad}

Para intentar hacer fallar el propio marco se definieron 25 clases de proceso entre los cinco benchmarks y cuatro niveles de evidencia:

\begin{description}
\item[T0:] geometría únicamente;
\item[T1:] geometría + microestructura + estructura + material;
\item[T2:] T1 + entorno + hidrología + térmica + acceso;
\item[T3:] T2 + duración + precisión + logística.
\end{description}

Cada proceso recibe

\begin{equation}
s(p,E)\in\{0,1,2\},
\end{equation}

donde 0 significa no recuperable, 1 favorecido pero no identificado y 2 recuperable a nivel de clase funcional bajo las reglas actuales.

Los resultados fueron:

\begin{table}[ht]
\centering
\begin{tabular}{lrrrr}
\toprule
Evidencia & Recuperables & Favorecidas & No recuperables & Score normalizado \\
\midrule
T0 & 1 & 0 & 24 & 4\% \\
T1 & 4 & 11 & 10 & 38\% \\
T2 & 10 & 9 & 6 & 58\% \\
T3 & 25 & 0 & 0 & 100\% \\
\bottomrule
\end{tabular}
\caption{Ablación determinista del formalismo actual. T3 no es inferencia terminal pura: incorpora objetivos y contexto constructivo.}
\end{table}

La conclusión es

\begin{equation}
\boxed{
X_T\ \text{aislado no identifica la mayoría de los procesos}.
}
\end{equation}

\section{Sensibilidad por modalidad}

En un análisis leave-one-modality-out desde T3, la caída de score fue:

\begin{center}
\begin{tabular}{lr}
\toprule
Modalidad eliminada & Caída de score \\
\midrule
Geometría & 32 \\
Microestructura & 22 \\
Logística & 13 \\
Estructura & 11 \\
Entorno & 9 \\
Material & 7 \\
Precisión & 7 \\
Schedule/duración & 5 \\
Accesibilidad & 4 \\
Térmica & 4 \\
Hidrología & 2 \\
\bottomrule
\end{tabular}
\end{center}

Estas cifras pertenecen al formalismo y conjunto de casos actuales; no son constantes universales.

Una ablación aleatoria exploratoria mostró la misma fragilidad multimodal: con probabilidad 25\% de eliminar cada modalidad, el número medio de clases plenamente recuperables cayó aproximadamente a 11 de 25.

\section{Falsos residuos de fuerza}

Una ontología demasiado gruesa puede transformar procesos modulares convencionales en aparentes fuerzas imposibles.

Escalas artificiales de peso monolítico:

\begin{align}
\text{torre Golden Gate} &\sim197\ {\rm MN},\\
\text{cable principal Golden Gate} &\sim107\ {\rm MN},\\
\text{estructura metálica Eiffel} &\sim71.6\ {\rm MN},\\
\text{acero estructural Empire State} &\sim563.7\ {\rm MN}.
\end{align}

Ninguna de estas magnitudes demuestra una interacción desconocida.

\begin{principle}[Residual ontológico]
\begin{equation}
\boxed{
\Fres\neq0
\not\Rightarrow
\text{nueva física}.
}
\end{equation}
Debe demostrarse que el residuo sobrevive a refinamiento ontológico, procesos temporales, mecanismos convencionales e incertidumbre.
\end{principle}

\section{Correcciones epistemológicas derivadas}

El estudio adversarial obliga a corregir varias formulaciones iniciales:

\begin{itemize}
\item Golden Gate: se recupera la clase acceso aéreo temporal bajo sitio/acceso; no el catwalk exacto.
\item Eiffel: la modularidad es fuerte; jacks, sandboxes y grúa concreta requieren información adicional.
\item Hoover: el control térmico es fuerte; cooling pipes exactas requieren además un schedule exigente.
\item Empire State: JIT, pipeline extremo y jumping derricks no son inferencias terminales puras.
\item Sydney: el generador geométrico común es fuerte; erection arch y fábrica concreta requieren estructura, acceso y logística.
\end{itemize}

La consecuencia metodológica es

\begin{equation}
\boxed{
\text{clase funcional recuperable}
\neq
\text{mecanismo histórico exacto}.
}
\end{equation}

\section{Protocolo operativo de Meta-Harness}

El orden recomendado es

\begin{equation}
\boxed{
\begin{aligned}
&\text{CHECK IDENTIFIABILITY}\\
&\rightarrow
\text{INFER CAUSAL ORDER}\\
&\rightarrow
\text{INFER TEMPORARY ONTOLOGY}\\
&\rightarrow
\text{EXHAUST CONVENTIONAL PROCESSES}\\
&\rightarrow
\text{COMPUTE RESIDUAL}.
\end{aligned}
}
\end{equation}

Cada claim debe adjuntar clase epistemológica, evidencia, incertidumbre, contradicciones, identificabilidad, falsificadores y procedencia del modelo.
